\section{Research Outline}

Hytale is a sandbox RPG whose identity is closely tied to modding, creator tools, and community-made content. In many games, modding begins after release, when players reverse-engineer systems, edit assets, or build unofficial tools to make the game do something new. Hytale changes this relation by presenting creation as part of the game's basic promise. Its website describes it as ``an adventure built for both creation and play'' \cite{hytaleOfficial}. This makes Hytale useful for asking what happens when modding stops being outside the game and becomes part of the development process itself. My report examines how Hytale turns modding into a space between play, community creativity, and game development.

Hytale mods are created through a mix of official and semi-official tools. The most important ones are the Hytale Asset Editor, direct data asset editing, Blockbench for models and animations, and Java server plugins for gameplay changes \cite{hytaleModdingStrategy}. These tools already show that Hytale modding is meant to include different kinds of creators: artists, builders, server owners, and programmers. Modding is not limited to changing textures or adding small cosmetic details. It can also involve new systems, multiplayer rules, custom items, and larger changes to how the game behaves.

The main public distribution platform for Hytale mods is CurseForge, especially because public submissions, collections, and official contests are organized there \cite{curseforgeHytale}. This matters because mods are not only private experiments on a local computer. They become public artifacts that can be shared, downloaded, updated, and discussed by a wider community. At the same time, Hytale modding is not yet a completely finished system. It is an evolving infrastructure that is being developed alongside the game itself, which can be seen in the changing API and still-developing tooling.

A key part of Hytale's modding strategy is that it is server-side first. This means that server hosts choose which mods are active, while players can join modded servers without manually installing external files \cite{hytaleModdingStrategy}. Modding therefore becomes part of the multiplayer structure of the game, not just an individual player modification. This is important for my research question because it changes the role of modding: it is not only something a single player adds to their own copy of the game, but something that can define a shared server experience.

Hytale modding involves several different types of creators because its main formats are split between packs and plugins. Packs are more focused on assets and content, such as blocks, mobs, items, textures, models, and cosmetics. Plugins are more code-based and can affect commands, events, server systems, minigames, and deeper gameplay rules. This creates a broad modding landscape. Some creators work visually, producing models, textures, decorative objects, or cosmetic content. Others work as plugin developers, designing new mechanics and changing how the game behaves. Server owners and world builders also play an important role, because they combine different mods into playable multiplayer spaces. Finally, toolmakers and documentation writers support the ecosystem by creating tutorials, templates, and guides. Hytale modding therefore includes artists, programmers, designers, server hosts, and community organizers rather than one single type of modder.

For significant examples, I do not want to focus only on one specific mod, but on what certain mods and modders reveal about Hytale's ecosystem. One important case is Violet, the creator of \textit{Violet's Wardrobe}, a cosmetic mod that adds craftable customization options \cite{violetsWardrobe}. The mod itself is relevant because it shows the importance of visual and cosmetic creation in Hytale. However, Violet is even more important as a modder: after creating early Hytale mods, she began working with the Hytale team on official cosmetics \cite{violetHired}. This shows how modding can become a visible form of creative work and even a path into the official development process. In this case, modding is not only fan activity. It becomes a kind of portfolio and a bridge between community creation and professional game development.

The modding landscape has also changed because Hytale itself has had an unusual development history. The game was announced years ago, acquired by Riot Games, cancelled in 2025, and later revived after the original creators reacquired it \cite{hytaleRevival}. Reporting on the revival emphasizes that Hytale's return was connected to its original creator-focused vision and modding support from day one. This makes modding part of the game's identity, not just an optional feature added after release. Hytale is therefore not only a game that allows mods; it is a game whose public image and future are strongly tied to the promise of community creation.

Early Hytale modding also shows how unpredictable creator communities can be. Shortly after launch, players created experimental mods that ran Windows 95, Minecraft, and even Hytale itself inside Hytale \cite{gamesradarHytaleMods}. These are not necessarily typical gameplay mods, but they show that strong modding tools can turn a game into a platform for unexpected experiments. For my research question, this is important because it shows that once developers give players powerful creation tools, they also give up some control over what the game can become.

Overall, Hytale is a useful case study because it shows modding becoming part of the development pipeline. Mods are created with official tools, shared through public platforms, used to build multiplayer spaces, and sometimes recognized by the developers themselves. This changes the relationship between developers and players. The modder is no longer only someone who changes a finished game from the outside. In Hytale, the modder can become a creator, tester, community organizer, and potentially even part of the professional development process.